\documentclass[12pt,a4paper]{article}

\usepackage[T1]{fontenc}
\usepackage[utf8]{inputenc}
\usepackage[slovak,english]{babel}
\usepackage{lmodern}
\usepackage{microtype}
\usepackage{geometry}
\usepackage{booktabs}
\usepackage{tabularx}
\usepackage{array}
\usepackage{longtable}
\usepackage{enumitem}
\usepackage{amsmath}
\usepackage{amssymb}
\usepackage{siunitx}
\usepackage{xcolor}
\usepackage{hyperref}
\usepackage{csquotes}
\usepackage[backend=biber,style=ieee,maxcitenames=2]{biblatex}

\geometry{margin=2.5cm}
\hypersetup{
  colorlinks=true,
  linkcolor=blue!50!black,
  citecolor=blue!50!black,
  urlcolor=blue!50!black,
  pdftitle={Zhoda pracovných pozícií a zručností},
  pdfauthor={Tím projektu}
}
\addbibresource{refs.bib}

\newcommand{\code}[1]{\texttt{\detokenize{#1}}}
\newcommand{\model}[1]{\texttt{\detokenize{#1}}}

\title{\textbf{Zhoda pracovných pozícií a zručností}\\
\large Semantické vyhľadávanie a hodnotenie typov zručností v úlohe TalentCLEF 2026 Task B}
\author{Tím projektu}
\date{25. október 2026}

\begin{document}
\selectlanguage{slovak}
\maketitle

\begin{abstract}
Táto práca sa zaoberá úlohou TalentCLEF 2026 Task B, ktorej cieľom je
vyhľadať zručnosti relevantné pre názov pracovnej pozície a zoradiť ich podľa
relevantnosti. Zároveň úloha rozlišuje medzi základnými zručnosťami
(\emph{core}) a kontextovými zručnosťami (\emph{contextual}). Navrhli sme
reprodukovateľný experimentálny pipeline založený na sentence-embedding
modeloch, kosínovej podobnosti, reprezentácii ESCO zručností pomocou aliasov
a kontrastívnom fine-tuningu vybraných modelov. Porovnali sme päť
predtrénovaných modelov a tri stratégie spracovania aliasov. Najlepší výsledok
podľa graded nDCG dosiahol model \model{TechWolf/JobBERT-v2} bez fine-tuningu,
pri priemerovaní embeddingov všetkých aliasov zručnosti
(\(nDCG=0{,}7381\)). Fine-tuning zlepšil všeobecné modely E5 a MPNet, zatiaľ
čo pri JobBERT zlepšil najmä presnosť v horných pozíciách. Podľa dostupných
qrels modely lepšie oddeľujú posudzované relevantné položky od
nepozitívnych alebo nehodnotených kandidátov než zoradia základné a
kontextové zručnosti.
\end{abstract}

\noindent\textbf{Kľúčové slová:} spracovanie prirodzeného jazyka, sentence
embeddings, ESCO, vyhľadávanie zručností, ranking, TalentCLEF, JobBERT

\section{Problém, motivácia a ciele}

Digitálne pracovné portály, nástroje na odporúčanie pracovných ponúk a systémy
na analýzu trhu práce potrebujú prepájať neštruktúrovaný text pracovných
ponúk s jednotným opisom profesií a zručností. Rovnaká zručnosť môže byť v
texte vyjadrená rôznymi synonymami a názov pracovnej pozície často neposkytuje
dostatočne podrobný zoznam požiadaviek. Ručné mapovanie pozícií na zručnosti je
časovo náročné, ťažko škálovateľné a môže byť nekonzistentné medzi anotátormi
alebo organizáciami.

TalentCLEF 2026 Task B formalizuje tento problém ako vyhľadávanie zručností
pre názov pracovnej pozície a klasifikáciu typu nájdenej zručnosti
\parencite{talentclef2026overviewb}. V kontexte úlohy je základná zručnosť
potrebná na výkon práce bez ohľadu na konkrétneho zamestnávateľa alebo
pracovné prostredie. Kontextová zručnosť závisí viac od organizácie, odvetvia
alebo projektu a má preto voliteľnejší charakter.

Motiváciou projektu je vytvoriť systém, ktorý:
\begin{enumerate}[label=\arabic*.]
  \item automaticky vyhľadá zručnosti relevantné pre názov pozície;
  \item využije normalizovaný slovník ESCO namiesto voľného generovania názvov;
  \item zoradí zručnosti tak, aby sa relevantné položky nachádzali na začiatku;
  \item vyhodnotí rozdiel medzi základnou a kontextovou relevanciou
  prostredníctvom graded rankingu;
  \item bude reprodukovateľný a vyhodnotený na oficiálnom development formáte.
\end{enumerate}

Projekt skúma tieto výskumné otázky:
\begin{description}
  \item[VO1:] Ako dobre dokážu predtrénované sentence-embedding modely
  vyhľadávať ESCO zručnosti pre názov pozície?
  \item[VO2:] Zlepšuje doménová špecializácia modelu výsledky oproti
  všeobecným embedding modelom?
  \item[VO3:] Pomáha pri reprezentácii zručnosti použiť všetky jej aliasy?
  \item[VO4:] Zlepšuje kontrastívny fine-tuning vybraných modelov retrieval?
  \item[VO5:] Sú modely úspešnejšie pri rozlíšení relevantnej a irelevantnej
  zručnosti než pri zoradení základných a kontextových zručností?
\end{description}

\section{Súvisiaca práca}

\subsection{TalentCLEF a ESCO}

TalentCLEF rozširuje klasické úlohy informačného vyhľadávania o doménu
ľudského kapitálu a pracovných zručností. Prehľad súťaže opisuje spoločné
úlohy, dáta a hodnotenie a zdôrazňuje potrebu normalizovaných znalostí o
pracovných pozíciách a zručnostiach \parencite{gasco2026overview}.
Špecifický prehľad Task B definuje párovanie názvov pozícií so zručnosťami a
graded relevanciu, v ktorej sa rozlišujú základné a kontextové zručnosti
\parencite{talentclef2026overviewb}.

ESCO predstavuje európsku viacjazyčnú klasifikáciu zručností, kompetencií,
kvalifikácií a povolaní. V tomto projekte sa využívajú ESCO identifikátory a
textové termíny zručností. Identifikátor umožňuje, aby sa synonymné alebo
alternatívne pomenovania spájali s jednou normalizovanou položkou.

\subsection{Prístupy účastníkov}

Predchádzajúce riešenia v rámci TalentCLEF kombinujú lexikálne, neurónové a
hybridné retrieval metódy. Práca MARSAD opisuje ESCO-aware hybridné
vyhľadávanie, v ktorom sa kombinuje informácia zo štruktúry ESCO a textová
podobnosť \parencite{IbrahimMARSADAT}. Tento smer je relevantný pre náš projekt,
pretože ukazuje, že samotná povrchová podobnosť nemusí zachytiť vzťahy medzi
pozíciami a zručnosťami.

Tím ui-nlp používa bi-encoder retrieval a ensemble stratégie pre úlohy
job-person a job-skill matching \parencite{LuthfiyyahuinlpAT}. Bi-encoder
prístup je blízky nášmu pipeline: dotazy a kandidátne dokumenty sa kódujú
samostatne a následne sa porovnávajú vo vektorovom priestore. Výhodou je
škálovateľnosť na tisíce kandidátnych zručností; nevýhodou je, že samostatné
embeddingy nemusia zachytiť všetky jemné interakcie medzi dvojicou textov.

Naše riešenie preto používa sentence embeddings ako transparentný a
reprodukovateľný baseline. Rozširuje ho porovnaním doménového modelu
JobBERT, všeobecných modelov a kontrastívneho fine-tuningu. Nejde o
generovanie nových zručností: všetky výsledky sa vyberajú z poskytnutého
ESCO-normalizovaného corpus.

\section{Dáta}

\subsection{Pôvod a formát}

Použité dáta pochádzajú z adresára \code{data/raw/TaskB}. Tréningová časť
obsahuje vzťahy medzi povolaniami a zručnosťami odvodené z ESCO a dve mapovania
identifikátorov na textové termíny. Development časť obsahuje dotazy, úplný
slovník kandidátnych zručností a súbor qrels s anotovanou relevanciou.

V tejto práci používame nasledujúce pojmy:
\begin{description}
  \item[Povolanie (occupation):] identifikované ESCO URI alebo tréningovým
  \code{job_id}; jeho textové pomenovanie je názov pozície.
  \item[Zručnosť (skill):] ESCO položka identifikovaná URI alebo \code{skill_id}.
  \item[Alias:] alternatívny textový názov tej istej ESCO zručnosti.
  \item[Query:] názov pracovnej pozície, pre ktorú sa vyhľadávajú zručnosti.
  \item[Corpus item:] kandidátna ESCO zručnosť, ktorú možno zaradiť do
  výsledného zoznamu.
  \item[Qrel:] judgment vzťahu query--corpus item. Obsahuje query ID,
  corpus ID a graded relevanciu.
\end{description}

\begin{table}[ht]
\centering
\caption{Počty záznamov v použitých dátových súboroch.}
\label{tab:data}
\begin{tabularx}{\textwidth}{>{\raggedright\arraybackslash}XrX}
\toprule
\textbf{Súbor} & \textbf{Počet} & \textbf{Význam} \\
\midrule
\code{training/job2skill.tsv} & 114\,699 & tréningové väzby povolanie--zručnosť \\
\code{training/jobid2terms.json} & 3\,039 & top-level identifikátory povolaní \\
\code{training/skillid2terms.json} & 13\,939 & top-level identifikátory zručností \\
\code{validation/queries} & 304 & development query názvov pozícií \\
\code{validation/corpus_elements} & 9\,052 & kandidátne ESCO zručnosti \\
\code{validation/qrels.tsv} & 56\,418 & anotované query--skill judgments \\
\code{processed/train.csv} & 164\,699 & páry po pridaní vzorkovaných negatívov \\
\code{processed/test.csv} & 96\,418 & lokálny validation-derived súbor \\
\bottomrule
\end{tabularx}
\end{table}

Počty v tabuľke predstavujú fyzické dátové záznamy. JSON súbory obsahujú
3\,039, respektíve 13\,939 top-level identifikátorov; po spojení s
tréningovým mappingom môže byť počet aktívne prepojených identifikátorov
menší, čo však nie je alternatívna definícia počtov uvedených v tabuľke.
Pri súboroch TSV \code{job2skill.tsv} a \code{qrels.tsv} nejde o riadok s
hlavičkou. Pri spracovaní sa k pozitívnym riadkom pridajú textové názvy a
následne sa do CSV doplnia syntetické negatívy. Preto má
\code{processed/train.csv} \(114\,699+50\,000=164\,699\) dátových riadkov
a \code{processed/test.csv} \(56\,418+40\,000=96\,418\) dátových riadkov;
CSV hlavička sa do týchto počtov nezapočítava. Qrels explicitne obsahuje
grade \(2\) pre core zručnosti a grade \(1\) pre contextual zručnosti.
Kandidátne skills, ktoré v qrels chýbajú, dostávajú pri evaluácii grade \(0\).
V qrels je 16\,161 záznamov grade \(2\) a 40\,257 záznamov grade \(1\).

\subsection{Licencia a redistribúcia}

V pracovnom priestore nie je pri dátových súboroch priložený samostatný
licenčný súbor, ktorý by udeľoval všeobecné právo na redistribúciu. Dáta preto
neinterpretujeme ako voľne redistribuovateľný dataset. Použitie ESCO údajov a
materiálov TalentCLEF musí rešpektovať podmienky organizátorov, podmienky
zdroja ESCO a prípadné podmienky konkrétneho kurzu alebo súťaže. Táto správa
opisuje lokálny experiment; pred zverejnením dát alebo modelových artefaktov
je potrebné overiť aktuálne licenčné a redistribučné pravidlá.

\section{Predspracovanie dát}

Implementácia predspracovania je v notebooku
\code{data_processing.ipynb}. Na začiatku sa načíta \code{job2skill.tsv},
mapovanie \code{jobid2terms.json} a mapovanie
\code{skillid2terms.json}. Z každého zoznamu termínov sa pre lokálny
spracovaný súbor vyberie prvý termín ako \code{job_name}, respektíve
\code{skill_name}. Následne sa obe mapovania spoja s tabuľkou väzieb podľa
identifikátorov.

Pôvodné označenia tréningových väzieb sú:
\begin{center}
\begin{tabular}{lll}
\toprule
\textbf{Pôvodné označenie} & \textbf{Číselný label} & \textbf{Interpretácia} \\
\midrule
\code{essential} & 2 & základná zručnosť \\
\code{optional} & 1 & kontextová/voliteľná zručnosť \\
\code{irrelevant} & 0 & irelevantná dvojica \\
\bottomrule
\end{tabular}
\end{center}

\subsection{Vzorkovanie negatívnych párov}

Tréningové mapovanie obsahuje iba známe pozitívne väzby, preto notebook
vytvára 50\,000 negatívnych párov. Náhodne sa vyberajú povolania a zručnosti
zo zoznamov identifikátorov. Dvojica sa prijme iba vtedy, ak sa nenachádza v
množine známych pozitívnych dvojíc. Proces opakuje bloky 1\,000 náhodných
výberov, až kým nevznikne 50\,000 unikátnych kandidátov. Seed je nastavený
na hodnotu \(42\) v súbore \code{config.py}. Páry sa následne spoja s textami
a dostanú label \(0\).

Takto vytvorené negatívy sú \emph{assumed negatives}: ide o dvojice, ktoré
nie sú v poskytnutom mappingu, nie o manuálne overené negatívne anotácie.
Neprítomnosť väzby preto nemusí znamenať, že zručnosť je v realite pre dané
povolanie nesprávna. Toto obmedzenie je dôležité pri interpretácii modelov,
ktoré by tieto negatívy používali ako učebné príklady.

Pre validation-derived \code{test.csv} sa rovnakým postupom pridá 40\,000
náhodných negatívnych párov k oficiálnym qrels z development časti. Tento
lokálny súbor nie je skrytý oficiálny test súťaže a nepoužíva sa ako hlavný
ranking benchmark. Hlavné experimenty pracujú priamo s 304 development
queries, úplným corpusom 9\,052 skills a qrels.

\section{Finálna metodika a experimentálny pipeline}

\subsection{Spoločný ranking pipeline}

Finálny experimentálny pipeline je implementovaný v tomto poradí:
\begin{enumerate}
  \item \code{data_processing.ipynb}: príprava tréningových a
  validation-derived tabuliek;
  \item \code{04_five_model_comparison.ipynb}: porovnanie piatich modelov
  a troch aliasových stratégií;
  \item \code{05_locked_model_comparison.ipynb}: uzamknutie mean
  reprezentácie a výber modelov na fine-tuning;
  \item \code{06_finetune_selected_models.ipynb}: fine-tuning JobBERT,
  E5 a MPNet;
  \item \code{07_final_comparison_and_error_analysis.ipynb}: finálne
  porovnanie a error analysis.
\end{enumerate}

Spoločný ranking postup v notebookoch 04--07 pozostáva z týchto krokov:
\begin{enumerate}
  \item načítanie development queries, corpus elements a qrels;
  \item zakódovanie názvov pozícií a textov skill aliasov sentence-embedding
  modelom;
  \item vytvorenie reprezentácie každej skill;
  \item výpočet kosínovej podobnosti medzi query embeddingom a každou
  kandidátnou skill;
  \item zoradenie všetkých 9\,052 kandidátov pre každú z 304 pozícií;
  \item porovnanie výsledného poradia s qrels.
\end{enumerate}

Embeddings sa normalizujú, takže kosínová podobnosť je v implementácii
ekvivalentná skalárnemu súčinu. Modely neprodukujú nové názvy skills; skóre
určuje iba poradie položiek v poskytnutom gazetteer.

\subsection{Porovnávané modely}

Notebook \code{04_five_model_comparison.ipynb} porovnáva:
\begin{itemize}
  \item \model{all-MiniLM-L6-v2};
  \item \model{all-mpnet-base-v2};
  \item \model{TechWolf/JobBERT-v2};
  \item \model{BAAI/bge-base-en-v1.5};
  \item \model{intfloat/e5-base-v2}.
\end{itemize}

Pre každú zručnosť sa testujú tri reprezentácie:
\begin{description}
  \item[first:] použije sa iba prvý alias;
  \item[mean:] embeddingy všetkých aliasov sa spriemerujú a výsledný vektor
  sa znovu normalizuje;
  \item[max:] skóre zručnosti je maximum kosínovej podobnosti cez všetky
  jej aliasy.
\end{description}

Notebook 05 následne \emph{uzamkne} globálnu stratégiu aliasov. Kritériom je
priemerný graded nDCG cez všetkých päť modelov; pri zhode sa použije MAP.
Výsledkom je stratégia \code{mean}: \(0{,}6526\) oproti \(0{,}6502\) pre
\code{first} a \(0{,}6426\) pre \code{max}. Na základe poradia modelov pri
uzamknutej reprezentácii sa na fine-tuning vyberú JobBERT, E5 a MPNet.

\paragraph{Rozsah implementovanej úlohy.}
Projekt implementuje retrieval s graded relevanciou, nie samostatný
trojtriedny klasifikátor. Grade \(2\) označuje core/essential skill, grade
\(1\) contextual/optional skill a kandidát bez záznamu v qrels dostáva pri
evaluácii grade \(0\). Fine-tuning používa labels \(1\) a \(2\) ako pozitívne
páry. Rozlíšenie core a contextual sa preto vyhodnocuje nepriamo poradím
prostredníctvom graded nDCG. Samostatná classification head, ktorá by
predikovala core, contextual, complementary alebo transversal triedu,
implementovaná nie je. Širšie termíny complementary a transversal sa v tomto
projekte nepoužívajú ako samostatné labely.

\subsection{Reprodukovateľnosť experimentálneho pipeline}

Projekt poskytuje zdokumentovaný a reprodukovateľný experimentálny postup.
Presná numerická reprodukcia však nemusí byť za každých okolností identická
v dôsledku verzií softvéru, použitého hardvéru, konfigurácie CUDA,
nondeterministického tréningu a opätovného sťahovania modelov. Lokálne
fine-tuned checkpointy nie sú súčasťou repozitára. Experimenty používajú
\code{RANDOM_SEED = 42}. Závislosti sú uvedené v
\code{requirements.txt}: \code{pandas}, \code{numpy}, \code{scikit-learn},
\code{sentence-transformers}, \code{jupyter}, \code{datasets},
\code{accelerate} a \code{matplotlib}; ich verzie však nie sú úplne
pripnuté. Predtrénované modely sa pri prvom spustení načítajú z Hugging Face.
Notebooky je potrebné spúšťať z koreňa repozitára, aby relatívne cesty k
dátam fungovali.

\section{Metriky a experimentálne nastavenie}

\subsection{Metriky}

Implementácia metrík je v súbore \code{ranking_metrics.py}. Pre každý query
sa najprv vytvorí stabilné poradie podľa klesajúceho skóre. Stabilné triedenie
zabezpečuje deterministické správanie pri zhodných skóre.

\paragraph{Graded nDCG.}
Hlavná metrika \emph{graded nDCG} používa zisky \(2\) pre základné a \(1\)
pre kontextové zručnosti. Diskontovanie podľa pozície znamená, že relevantná
zručnosť na začiatku zoznamu je cennejšia než rovnaká zručnosť na jeho konci:
\[
DCG = \sum_{i=1}^{n}\frac{rel_i}{\log_2(i+1)},\qquad
nDCG = \frac{DCG}{IDCG}.
\]
Táto metrika je vhodná, pretože zohľadňuje retrieval aj rozdiel medzi
core a contextual.

\paragraph{Binary nDCG a MAP.}
Pri binárnych metrikách sa labels \(1\) a \(2\) zlúčia do hodnoty relevantné.
MAP meria, ako kompaktne sú všetky relevantné položky umiestnené v poradí,
nie iba prvá nájdená položka. Binary nDCG a MAP preto hodnotia schopnosť
odfiltrovať irelevantné zručnosti.

\paragraph{MRR a Precision@K.}
MRR používa recipročnú pozíciu prvej relevantnej zručnosti. Precision@1,
Precision@5 a Precision@10 merajú priemerný podiel relevantných zručností
v horných \(K\) pozíciách. Sú praktické pre používateľské rozhrania, ktoré
zobrazujú krátky zoznam odporúčaní.

Žiadna z týchto metrík sama osebe nie je klasická presnosť
trojtriedneho klasifikátora. V tomto projekte model produkuje rankingové
skóre a graded nDCG zachytáva rozdiel core--contextual nepriamo cez poradie.

\subsection{Experimentálny protokol}

Po výbere mean alias reprezentácie sa v notebooku
\code{06_finetune_selected_models.ipynb} fine-tunujú tri vybrané modely.
Použijú sa iba pozitívne páry s \code{label > 0}; labels \(1\) a \(2\) sú
teda počas fine-tuningu považované za rovnocenné pozitívne páry. Nastavenia
experimentu sú:
\begin{center}
\begin{tabular}{ll}
\toprule
Loss & \code{MultipleNegativesRankingLoss} \\
Počet epoch & 2 \\
Batch size & 128 \\
Learning rate & \(2\cdot10^{-5}\) \\
Seed & 42 \\
\bottomrule
\end{tabular}
\end{center}

Pri tomto loss je správna skill partnerom názvu pozície a ostatné skills z
batchu slúžia ako in-batch negatives. Sampler
\code{NO_DUPLICATES} odstraňuje presné opakovanie textov v batchi, ale
neodstraňuje sémanticky podobné false negatives. Keďže jedna skill môže byť
relevantná pre viac súvisiacich pozícií, ide o dôležité obmedzenie.
Tréningový dataset používa stĺpce \code{anchor} a \code{positive}; tieto
stĺpce sú v konfigurácii Sentence Transformers odovzdané priamo do
zodpovedajúcich rout modelu JobBERT. 

Rozdelenie internej validácie sa robí na úrovni názvu pozície, nie náhodným
rozdelením jednotlivých riadkov. Tak sa znižuje riziko, že sa rovnaká
pozícia objaví v tréningu aj internej validácii. Finálna evaluácia sa však
vykonáva jednotným protokolom na celom development sete v notebooku
\code{07_final_comparison_and_error_analysis.ipynb}.

Notebook 07 porovnáva osem systémov: päť zero-shot systémov a tri
fine-tunované systémy. Okrem agregovaných metrík vytvára výsledky po
jednotlivých pozíciách pre všetkých 304 queries, súbor s najťažšími
pozíciami a AUC analýzu rozlíšenia relevantných položiek a core položiek.

\section{Výsledky}

\begin{table}[ht]
\centering
\caption{Hlavné výsledky na development sete; vyššia hodnota je lepšia.}
\label{tab:results}
\begin{tabular}{llrrrrrr}
\toprule
\textbf{Fáza} & \textbf{Model} & \textbf{nDCG} & \textbf{MAP} &
\textbf{MRR} & \textbf{P@1} & \textbf{P@5} & \textbf{P@10}\\
\midrule
Zero-shot & JobBERT & 0.7381 & 0.3416 & 0.7886 & 0.6480 & 0.5625 & 0.5178\\
Fine-tuned & JobBERT & 0.7319 & 0.3258 & \textbf{0.7997} & \textbf{0.6678} & \textbf{0.5882} & \textbf{0.5507}\\
Fine-tuned & E5 & 0.7040 & 0.2694 & 0.7243 & 0.5658 & 0.5138 & 0.5039\\
Fine-tuned & MPNet & 0.6981 & 0.2602 & 0.7071 & 0.5625 & 0.5033 & 0.4766\\
Zero-shot & E5 & 0.6401 & 0.1642 & 0.6958 & 0.5428 & 0.4507 & 0.4168\\
Zero-shot & MPNet & 0.6327 & 0.1612 & 0.6302 & 0.4539 & 0.3980 & 0.3711\\
Zero-shot & BGE & 0.6309 & 0.1555 & 0.6340 & 0.4441 & 0.4138 & 0.3737\\
Zero-shot & MiniLM & 0.6213 & 0.1476 & 0.6348 & 0.4704 & 0.3836 & 0.3655\\
\bottomrule
\end{tabular}
\end{table}

Ako autoritatívny zdroj agregovaných výsledkov používame súbor
\code{results/final_system_comparison.csv}, ktorý vznikol vo finálnom
porovnávacom notebooku. V tabuľke je preto pre zero-shot JobBERT uvedené
\(MRR=0{,}7886\). Samostatný export error analysis dopočítava metriky z
per-job tabuliek a uvádza približne \(0{,}7903\); tento rozdiel je dôsledkom
odlišného zdroja agregácie, nie zmeny modelu. Fine-tuned JobBERT má v oboch
výstupoch \(MRR=0{,}7997\). Pre porovnateľnosť sú všetky čísla v tabuľke
prebraté z jedného autoritatívneho CSV.

JobBERT bez fine-tuningu dosahuje najlepší graded nDCG a MAP. Tento výsledok
je konzistentný s možným prínosom doménovej špecializácie v tejto úlohe,
nepredstavuje však test štatistickej významnosti. Fine-tuning E5 zvyšuje
nDCG z \(0{,}6401\)
na \(0{,}7040\) a MPNet z \(0{,}6327\) na \(0{,}6981\). Pri JobBERT sa
celkový graded nDCG mierne zníži, ale zlepší sa MRR a všetky sledované
Precision@K. To znamená, že fine-tuned JobBERT častejšie umiestni relevantnú
zručnosť úplne na začiatok, hoci celé graded poradie je podľa nDCG o niečo
slabšie. Podľa tejto konfigurácie teda fine-tuning zlepšil reportované
retrieval metriky pre E5 a MPNet a top-k metriky pre JobBERT, zatiaľ čo
zero-shot JobBERT zostal najsilnejší podľa graded nDCG. Štatistické testy
významnosti neboli vykonané.

Notebook 07 uvádza aj AUC výsledky. Zero-shot JobBERT dosahuje približne
\(0{,}9596\) pre rozlíšenie relevantných a irelevantných zručností a
\(0{,}9585\) pre preferenciu core voči contextual. Fine-tuned JobBERT
dosahuje \(0{,}9420\), respektíve \(0{,}9483\). Tieto výsledky podporujú
hlavný záver pre dostupné qrels: modely dobre oddeľujú posudzované relevantné
skills od nepozitívnych alebo nehodnotených kandidátov. Tieto hodnoty však
nie sú dôkazom, že každý kandidát bez qrels je skutočne irelevantný, a AUC
nie je presnosť trojtriedneho klasifikátora. Rozdiel medzi core a contextual
je jemnejší a v tomto projekte sa nepredikuje samostatnou classification head.

\section{Možné rozšírenia}

Nasledujúce návrhy nie sú súčasťou implementovaného pipeline:
\begin{itemize}
  \item graded alebo pairwise ranking loss, ktorý by počas tréningu
  rozlišoval grade \(2\) od grade \(1\);
  \item filtrovanie false negatives pri in-batch negatívoch;
  \item explicitný klasifikátor pre core a contextual labels;
  \item kvalitnejšie alebo manuálne overené negatívne vzorky;
  \item intervaly spoľahlivosti a párové testy štatistickej významnosti.
\end{itemize}

\section{Použitie AI a reprodukovateľnosť}

\subsection{Použitie LLM a AI agentov}

Počas projektu boli LLM použité ako podporný nástroj pri návrhu pipeline,
písaní prvých verzií kódu, interpretácii chýb a kontrole dokumentácie.
Experimenty a výsledky však vznikli lokálnym spustením notebookov členmi
tímu. LLM nebola zdrojom nameraných metrík. Vývojový záznam v súbore
\code{llm.md} uvádza použité nástroje, etapy práce, overovanie výstupov aj
zistené chyby. Každý návrh vygenerovaný LLM bol preto potrebné skontrolovať
voči kódu, dátam a skutočnému výstupu experimentu.

\subsection{Postup reprodukcie}

Reprodukciu možno vykonať nasledovne:
\begin{enumerate}
  \item naklonovať \href{https://github.com/TheDim0nu4/Job-Skill-Matching-Ranking-and-Classification-Using-NLP-Approaches}{\underline{repozitár}} a spustiť terminál v jeho koreňovom adresári;
  \item vytvoriť kompatibilné Python prostredie;
  \item nainštalovať závislosti príkazom
  \code{pip install -r requirements.txt};
  \item skontrolovať, že \code{config.py} obsahuje
  \code{RANDOM_SEED = 42};
  \item spustiť \code{data_processing.ipynb}, ktorý vytvorí
  \code{data/processed/train.csv} a \code{data/processed/test.csv};
  \item spustiť \code{04_five_model_comparison.ipynb}, ktorý vytvorí
  porovnanie 15 konfigurácií;
  \item spustiť \code{05_locked_model_comparison.ipynb}, ktorý vyberie
  globálnu reprezentáciu aliasov a kandidátov na fine-tuning;
  \item spustiť \code{06_finetune_selected_models.ipynb}; notebook stiahne
  základné modely a uloží lokálne fine-tuned váhy;
  \item spustiť \code{07_final_comparison_and_error_analysis.ipynb}, ktorý
  vytvorí finálne tabuľky, per-job analýzu a export najťažších pozícií.
\end{enumerate}

Notebooky 04--07 tvoria finálnu metodiku tejto práce. Staršie notebooky
01--03 zostávajú v repozitári ako vývojová história, ale nie sú súčasťou
definície finálneho experimentálneho protokolu. Výsledky sa ukladajú do
adresára \code{results/}; kontrolnými výstupmi sú najmä
\code{final_system_comparison.csv}, \code{error_analysis_summary.csv} a
\code{error_analysis_per_job.csv}.

\section{Záver}

Navrhli a experimentálne sme overili zdokumentovaný experimentálny postup pre
vyhľadávanie ESCO zručností podľa názvu pracovnej pozície. Pipeline porovnáva
päť sentence-embedding modelov, tri stratégie reprezentácie aliasov a
kontrastívny fine-tuning troch vybraných modelov. Hlavným výsledkom je, že
doménový JobBERT bez fine-tuningu dosahuje najlepší celkový graded ranking
(\(nDCG=0{,}7381\)), zatiaľ čo fine-tuned JobBERT poskytuje najlepšie
top-k výsledky.

Fine-tuning všeobecných modelov E5 a MPNet prináša zlepšenie v použitej
konfigurácii, zatiaľ čo pri JobBERT zvyšuje najmä top-k metriky. Súčasne
výsledky ukazujú, že retrieval a klasifikácia typu zručnosti nie sú totožné
problémy. Systém pracuje s graded rankingom a neobsahuje samostatný
klasifikátor pre core a contextual labels.

\printbibliography[heading=bibintoc,title={Literatúra}]

\end{document}
